% !PRINTED 8
\weberpage{8}

Auch wir wollen in der Folge nicht auf das Hülfsmittel der
geometrischen Anschauung verzichten, und z.~B. die Punkte einer
geraden Linie unbedenklich als eine stetige Menge betrachten.

\begin{center}
\rule{0.18\linewidth}{0.4pt}
\end{center}

Eine geordnete Menge $\mathfrak M$ heisst \gesperrt{messbar} unter
folgenden Voraussetzungen: Addition und Vervielfältigung sind in
$\mathfrak M$ allgemein ausführbar, ebenso Subtraction eines kleineren
von einem grösseren Element, d.~h. aus irgend zwei Elementen $a,b$
(die auch identisch sein können) kann nach einer bestimmten Vorschrift
ein neues Element $a+b$ von $\mathfrak M$ abgeleitet werden, so dass
$a+b$ grösser als $a$ und als $b$ ist, und dass die bekannten, in den
Formeln
\[
 a+b=b+a,
 \qquad (a+b)+c=a+(b+c)
\]
ausgedrückten Regeln der Addition gelten; und zu zwei Elementen $a,c$,
von denen das zweite grösser ist, kann ein drittes Element $b$
gefunden werden, so dass $a+b=c$ ist, was auch durch das Zeichen der
Subtraction $b=c-a$ ausgedrückt wird. Zwei ungleiche Elemente von
$\mathfrak M$ haben also immer eine bestimmte Differenz. Aus diesen
Voraussetzungen folgt, dass eine Summe grösser wird, wenn einer der
Summanden sich vergrössert. Die wiederholte, etwa $m$-malige Addition
desselben Elementes $a$ heisst Vervielfältigung und ihr Ergebniss wird
mit $ma$ bezeichnet.

Es kommt endlich noch eine Voraussetzung hinzu, nämlich die, dass bei
jedem gegebenen $a$ ein hinlänglich hohes Vielfaches $ma$ grösser ist
als ein beliebig gegebenes anderes Element $b$. Unter den Elementen
einer messbaren Menge giebt es also kein grösstes.

In einer dichten messbaren Menge giebt es auch kein kleinstes Element;
denn wäre $a$ das kleinste und $b$ ein beliebiges Element, so könnten
zwischen $b$ und $b+a$ keine Elemente liegen, weil, wenn $b<c<b+a$
wäre, aus der Definition der Messbarkeit folgen würde, dass $a'=c-b$
kleiner als $a$ wäre. Es folgt auch umgekehrt, dass eine messbare
Menge, in der kein kleinstes Element vorkommt, dicht ist. Denn sind
$a$ und $a+b$ irgend zwei Elemente, so braucht man ja zu $a$ nur ein
Element zu addiren, was kleiner als $b$ ist, um ein Element zwischen
$a$ und $a+b$ zu erhalten.

Ist eine Menge stetig, so lassen sich die Voraussetzungen für die
Messbarkeit noch vereinfachen, weil dann die Subtraction eine Folge
der Addition ist. Sind nämlich $a$ und $c$ zwei Ele-
%CONT>>
